\documentclass[11pt]{article}
\usepackage[margin=1in]{geometry}
\usepackage{enumitem}
% =============================================================================
% Preambles/header.tex — shared LaTeX/Beamer preamble for this project
%
% Use from a Beamer lecture by prepending:
%     \input{header}
% and compiling with TEXINPUTS=../Preambles:$TEXINPUTS (the /compile-latex
% skill does this automatically).
%
% PALETTE CONTRACT. Color names defined here MUST match the names in
% Quarto/theme-template.scss so Beamer and Quarto renderings use the same
% palette. The scripts/check-palette-sync.sh script enforces this.
%
% To customize: edit the HEX values below AND the matching $variable in
% Quarto/theme-template.scss, then run scripts/check-palette-sync.sh.
% =============================================================================

% -----------------------------------------------------------------------------
% Engine detection + encoding
%
% Our /compile-latex skill uses XeLaTeX, where inputenc is unnecessary and
% can produce encoding warnings. Only load inputenc under pdfTeX.
% -----------------------------------------------------------------------------
\usepackage{iftex}
\ifPDFTeX
  \usepackage[utf8]{inputenc}
\fi

\usepackage{amsmath, amssymb, amsthm}
\usepackage{booktabs}
\usepackage{graphicx}

% -----------------------------------------------------------------------------
% PALETTE — mirrors Quarto/theme-template.scss variable names.
% Load xcolor and define colors BEFORE anything that references them
% (notably hyperref, hypersetup, and the Beamer color assignments).
% -----------------------------------------------------------------------------
\usepackage{xcolor}

\definecolor{primary-blue}{HTML}{012169}      % headings, accents
\definecolor{primary-gold}{HTML}{B9975B}      % emphasis, borders
\definecolor{highlight-yellow}{HTML}{F2A900}  % markers, alerts
\definecolor{light-bg}{HTML}{E8EDF5}          % title / section backgrounds
\definecolor{jet}{HTML}{1A1A1A}               % body text

% Semantic colors (used in diagrams and annotations)
\definecolor{positive}{HTML}{15803D}          % good / observed / identified
\definecolor{negative}{HTML}{B91C1C}          % bad / problematic / bias
\definecolor{neutral}{HTML}{525252}           % reference / context

% Highlight accents (match .hi-* classes in the SCSS)
\definecolor{hi-slate}{HTML}{314F4F}
\definecolor{hi-green}{HTML}{2E8B57}
\definecolor{hi-red}{HTML}{C41E3A}

% Semantic aliases so skill templates and snippets can write
% \textcolor{accent}{...} without hard-coding "primary-blue".
\colorlet{accent}{primary-blue}
\colorlet{accent2}{primary-gold}

% -----------------------------------------------------------------------------
% Hyperref — loaded AFTER xcolor + palette so link colors resolve.
% -----------------------------------------------------------------------------
\usepackage{hyperref}
\hypersetup{colorlinks=true, linkcolor=primary-blue, urlcolor=primary-blue,
            citecolor=primary-blue, pdfborder={0 0 0}}

% -----------------------------------------------------------------------------
% Course code — set once per deck (after \input{header}, before \begin{document})
% via \coursecode{CS401}. Shown in the footer on every slide so a deck's course
% is identifiable without opening the title page. Empty by default.
% -----------------------------------------------------------------------------
\makeatletter
\newcommand{\coursecode}[1]{\gdef\@coursecodetext{#1}}
\gdef\@coursecodetext{}
\makeatother

% -----------------------------------------------------------------------------
% Beamer theme assignments (only apply under Beamer).
%
% We detect Beamer via \@ifclassloaded{beamer}, which is the documented,
% non-internal way. This must be wrapped in \makeatletter/\makeatother
% because \@ifclassloaded contains an @-catcode character.
% -----------------------------------------------------------------------------
\makeatletter
\@ifclassloaded{beamer}{%
  \usefonttheme{professionalfonts}
  \setbeamercolor{structure}{fg=primary-blue}
  \setbeamercolor{frametitle}{fg=primary-blue}
  \setbeamercolor{title}{fg=primary-blue}
  \setbeamercolor{subtitle}{fg=primary-gold}
  \setbeamercolor{author}{fg=primary-blue}
  \setbeamercolor{institute}{fg=neutral}
  \setbeamercolor{date}{fg=primary-gold}
  \setbeamercolor{itemize item}{fg=primary-blue}
  \setbeamercolor{itemize subitem}{fg=primary-gold}
  \setbeamercolor{alerted text}{fg=primary-gold}
  \setbeamercolor{block title}{fg=white, bg=primary-blue}
  \setbeamercolor{block body}{bg=light-bg}
  % Semantic block variants: block=definition/concept (blue), exampleblock=
  % worked example (green/positive), alertblock=key takeaway or caution (gold).
  % Keep to <=2 colored boxes per slide across ALL three variants (INV-7).
  \setbeamercolor{block title example}{fg=white, bg=positive}
  \setbeamercolor{block body example}{bg=light-bg}
  \setbeamercolor{block title alerted}{fg=white, bg=primary-gold}
  \setbeamercolor{block body alerted}{bg=light-bg}

  % Minimal footer: course code (left) + page X / N (right), no navigation clutter
  \setbeamertemplate{navigation symbols}{}
  \setbeamertemplate{footline}{%
    \color{neutral}\footnotesize\hspace{0.7em}\@coursecodetext\hfill
    \insertframenumber/\inserttotalframenumber\hspace{0.7em}\vspace{0.3em}}
}{}
\makeatother

% -----------------------------------------------------------------------------
% TikZ setup — libraries the snippet gallery uses
% -----------------------------------------------------------------------------
\usepackage{tikz}
\usetikzlibrary{arrows.meta, positioning, calc, decorations.pathreplacing,
                fit, shapes.geometric, backgrounds}

% Shared TikZ styles — snippets and hand-written diagrams can pick these up
% via `\begin{tikzpicture}[every node/.style={...}]` or by naming specific
% styles (e.g., `\node[dag-node] (X) at (0,0) {X};`).
\tikzset{
  % Node styles for causal DAGs and similar
  dag-node/.style = {
    draw=primary-blue, fill=light-bg, rounded corners=2pt,
    minimum width=1.2cm, minimum height=0.9cm, align=center, font=\small
  },
  decision-node/.style = {
    draw=primary-gold, fill=white, diamond, aspect=1.8,
    minimum width=1.8cm, minimum height=1.2cm, align=center, font=\small,
    inner sep=1pt
  },
  % Edge styles — observed vs counterfactual
  observed-edge/.style = {-{Latex[length=2.5mm]}, thick, draw=jet},
  counterfactual-edge/.style = {-{Latex[length=2.5mm]}, thick, draw=neutral, dashed},
  confound-edge/.style = {-{Latex[length=2.5mm]}, thick, draw=negative, dashed},
  % Data-point styles for plots
  observed-dot/.style = {fill=primary-blue, draw=primary-blue, circle, inner sep=1.2pt},
  counterfactual-dot/.style = {fill=white, draw=neutral, circle, inner sep=1.2pt}
}

% -----------------------------------------------------------------------------
% Convenience macros
% -----------------------------------------------------------------------------
\newcommand{\muted}[1]{\textcolor{neutral}{#1}}
\newcommand{\key}[1]{\textcolor{primary-gold}{\textbf{#1}}}
\newcommand{\good}[1]{\textcolor{positive}{#1}}
\newcommand{\bad}[1]{\textcolor{negative}{#1}}

% Standout frame macro: full-bleed dark background for section transitions.
% Beamer-only (uses \begin{frame}/\setbeamercolor/\usebeamerfont) -- do not
% call from an article-class file (e.g. Notes/<CODE>/*-notes.tex). \muted,
% \key, \good, \bad, and \coursecode above are plain xcolor/gdef and are
% safe to use from either class; verified when Notes/article-class support
% was added.
% Expand: \transitionslide{Your transition title}
\newcommand{\transitionslide}[1]{%
  \begingroup
  \setbeamercolor{background canvas}{bg=primary-blue}
  \begin{frame}[plain]
    \centering
    \vfill
    {\usebeamerfont{title}\color{white}\Large #1\par}
    \vfill
  \end{frame}
  \endgroup
}

% Section divider frame (Beamer-only), an alternative to \transitionslide with a
% darker two-line style: near-black (jet) background, a small white label line
% above a larger gold title, and a thin gold rule along the bottom edge.
% Expand: \sectiondivider{Section 2}{Pointers}
\newcommand{\sectiondivider}[2]{%
  \begingroup
  \setbeamercolor{background canvas}{bg=jet}
  \begin{frame}[plain]
    \vspace*{3.0cm}
    {\color{white}\ttfamily\large #1\par}
    \vspace{0.5cm}
    {\color{primary-gold}\LARGE #2\par}
    \begin{tikzpicture}[remember picture, overlay]
      \draw[primary-gold, line width=2.5pt]
        ($(current page.south west)+(0,0.1)$) -- ($(current page.south east)+(0,0.1)$);
    \end{tikzpicture}
  \end{frame}
  \endgroup
}

% -----------------------------------------------------------------------------
% Channel watermark — "Concepts That Click" (YouTube).
%
% Article-class files (CompetitiveExam Q/A sets, Notes, Assignments) get a
% light diagonal draft watermark on every page plus a centered footer naming
% the channel. Beamer decks get a subtle TikZ background watermark instead
% (the footline already carries the course code). Centralized here so every
% MCQ PDF is branded without per-file edits.
% -----------------------------------------------------------------------------
\makeatletter
\@ifclassloaded{article}{%
  \usepackage{draftwatermark}
  \SetWatermarkText{Concepts That Click}
  \SetWatermarkScale{1}
  \SetWatermarkFontSize{2cm}
  \SetWatermarkLightness{0.86}
  \SetWatermarkAngle{45}
  \usepackage{fancyhdr}
  \pagestyle{fancy}
  \fancyhf{}
  \fancyfoot[C]{\footnotesize\textcolor{neutral}{Concepts That Click~|~YouTube: \href{https://www.youtube.com/@concepts.thatclick}{@concepts.thatclick}}}
  \renewcommand{\headrulewidth}{0pt}
  \renewcommand{\footrulewidth}{0pt}
  \fancypagestyle{plain}{%
    \fancyhf{}%
    \fancyfoot[C]{\footnotesize\textcolor{neutral}{Concepts That Click~|~YouTube: \href{https://www.youtube.com/@concepts.thatclick}{@concepts.thatclick}}}%
    \renewcommand{\headrulewidth}{0pt}%
    \renewcommand{\footrulewidth}{0pt}%
  }
}{}
\@ifclassloaded{beamer}{%
  \setbeamertemplate{background}{%
    \begin{tikzpicture}[remember picture, overlay]
      \node[rotate=45, opacity=0.055, align=center]
        at (current page.center)
        {\Huge\textcolor{neutral}{Concepts That Click}};
    \end{tikzpicture}%
  }
}{}
\makeatother 
\coursecode{JKSSB}
\title{Lesson 30 Practice: Access Queries, Forms and Reports}
\author{JKSSB Computer Awareness}
\date{\today}
\begin{document}
\maketitle
\noindent\textbf{Provenance:} One item reproduces the concept of a real
Website Operator 2026 item (PYQ-39, provisional key); the remaining nineteen
are original JKSSB-pattern questions (\textit{[Original]}).

\begin{enumerate}[label=\textbf{Q\arabic*.}, leftmargin=*]
\item \textit{[Original]} The abbreviation \textbf{SQL} stands for:
  \begin{enumerate}[label=\Alph*.]
    \item Standard Query Language
    \item Structured Query Language
    \item Simple Query Language
    \item System Query Language
  \end{enumerate}
\item \textit{[Original]} Which MS Access object stores the actual data in
  rows and columns?
  \begin{enumerate}[label=\Alph*.]
    \item Query
    \item Form
    \item Table
    \item Report
  \end{enumerate}
\item \textit{[Original]} Which MS Access object is used to ask a question
  and retrieve selected data from the tables?
  \begin{enumerate}[label=\Alph*.]
    \item Form
    \item Query
    \item Report
    \item Table
  \end{enumerate}
\item \textit{[PYQ WO 2026 Q79, key:provisional]} A query that retrieves
  data on the basis of a condition is called a:
  \begin{enumerate}[label=\Alph*.]
    \item Delete query
    \item Update query
    \item Append query
    \item Select query
  \end{enumerate}
\item \textit{[Original]} Which type of query changes the existing values in
  records?
  \begin{enumerate}[label=\Alph*.]
    \item Select query
    \item Update query
    \item Append query
    \item Crosstab query
  \end{enumerate}
\item \textit{[Original]} Which type of query adds records from one table to
  another?
  \begin{enumerate}[label=\Alph*.]
    \item Update query
    \item Delete query
    \item Append query
    \item Select query
  \end{enumerate}
\item \textit{[Original]} Which type of query removes records that match a
  condition?
  \begin{enumerate}[label=\Alph*.]
    \item Delete query
    \item Select query
    \item Update query
    \item Append query
  \end{enumerate}
\item \textit{[Original]} Which of the following does \textbf{NOT} modify the
  stored data?
  \begin{enumerate}[label=\Alph*.]
    \item Update query
    \item Append query
    \item Delete query
    \item Select query
  \end{enumerate}
\item \textit{[Original]} Which MS Access object provides an on-screen
  interface for entering and viewing data one record at a time?
  \begin{enumerate}[label=\Alph*.]
    \item Table
    \item Query
    \item Form
    \item Report
  \end{enumerate}
\item \textit{[Original]} Which MS Access object presents data in a
  formatted, printable layout?
  \begin{enumerate}[label=\Alph*.]
    \item Form
    \item Report
    \item Query
    \item Table
  \end{enumerate}
\item \textit{[Original]} Which statement about an Access report is correct?
  \begin{enumerate}[label=\Alph*.]
    \item A report is used to enter new records.
    \item A report is read-only and is meant for printing.
    \item A report stores the actual data of the database.
    \item A report modifies records that match a condition.
  \end{enumerate}
\item \textit{[Original]} Which option gives the correct order of Access
  objects in which data is stored, selected, entered and printed?
  \begin{enumerate}[label=\Alph*.]
    \item Query $\rightarrow$ Table $\rightarrow$ Report $\rightarrow$ Form
    \item Table $\rightarrow$ Form $\rightarrow$ Query $\rightarrow$ Report
    \item Table $\rightarrow$ Query $\rightarrow$ Form $\rightarrow$ Report
    \item Form $\rightarrow$ Table $\rightarrow$ Query $\rightarrow$ Report
  \end{enumerate}
\item \textit{[Original]} Forms and reports in Access are usually built on:
  \begin{enumerate}[label=\Alph*.]
    \item another form.
    \item a table or a query.
    \item a report only.
    \item the Navigation Pane.
  \end{enumerate}
\item \textit{[Original]} In Access, a lookup field is used to:
  \begin{enumerate}[label=\Alph*.]
    \item delete records that match a condition.
    \item print a formatted summary of the data.
    \item choose a value from a list supplied by a table or query.
    \item create a new table automatically.
  \end{enumerate}
\item \textit{[Original]} Which of the following is \textbf{NOT} an object in
  an MS Access database?
  \begin{enumerate}[label=\Alph*.]
    \item Table
    \item Query
    \item Report
    \item Spreadsheet
  \end{enumerate}
\item \textit{[Original]} Evaluate the statements:
  \begin{enumerate}[label=\Roman*.]
    \item A Select query modifies the stored data.
    \item An Update query changes existing records.
    \item A Form is used to print a formatted report.
  \end{enumerate}
  \begin{enumerate}[label=\Alph*.]
    \item I and II only
    \item II only
    \item I and III only
    \item I, II and III
  \end{enumerate}
\item \textit{[Original]} The grid used in Access Query Design view to build
  a query by choosing fields and criteria is called:
  \begin{enumerate}[label=\Alph*.]
    \item QBE (Query By Example)
    \item SQL
    \item QWERTY
    \item PDF
  \end{enumerate}
\item \textit{[Original]} Which SQL keyword is used to retrieve data from a
  table?
  \begin{enumerate}[label=\Alph*.]
    \item UPDATE
    \item SELECT
    \item DELETE
    \item INSERT
  \end{enumerate}
\item \textit{[Original]} In the SQL statement
  ``\texttt{SELECT Name FROM Students WHERE Marks > 60}'', the
  \texttt{WHERE} clause:
  \begin{enumerate}[label=\Alph*.]
    \item names the table to be searched.
    \item lists the columns to be displayed.
    \item gives the condition that the rows must satisfy.
    \item sorts the result in ascending order.
  \end{enumerate}
\item \textit{[Original]} Which pairing is correctly matched?
  \begin{enumerate}[label=\Alph*.]
    \item Table --- asks a question and retrieves data
    \item Query --- stores the actual data
    \item Form --- provides a screen for data entry
    \item Report --- modifies records that match a condition
  \end{enumerate}
\end{enumerate}
\end{document}
